\chapter{Liver Infection}
\index{Liver Infection}

\section*{Definition and Overview}
A liver infection is a clinical condition caused by the invasion and multiplication of pathogenic microorganisms within the hepatic parenchyma or biliary system. The most common form is viral hepatitis, which involves inflammation of the liver caused by one of several hepatotropic viruses (primarily hepatitis A, B, C, D, and E). In addition to viral causes, liver infections can present as localized suppurative processes, such as pyogenic liver abscesses (caused by bacterial pathogens) or amoebic liver abscesses (caused by parasitic protozoa). These infections disrupt essential hepatic functions, including metabolic regulation, detoxification, and protein synthesis. While some liver infections are acute and self-limiting, others can establish chronic infection, leading to progressive liver injury, fibrosis, cirrhosis, and an increased risk of primary liver cancer.

\section*{Epidemiology}
Liver infections represent a significant global and domestic health burden. Viral hepatitis is the leading cause, with hepatitis B (HBV) and hepatitis C (HCV) causing chronic infections in hundreds of millions of people worldwide. In many countries, chronic hepatitis B affects approximately 200,000 people, with a higher prevalence among Aboriginal and Torres Strait Islander peoples and individuals born in regions with high endemicity (such as Asia-Pacific and Sub-Saharan Africa). Chronic hepatitis C affects around 110,000 people, primarily linked to past or current unsafe injecting drug use, though the introduction of highly effective direct-acting antiviral (DAA) therapies has led to a rapid decline in prevalence. Pyogenic liver abscesses are relatively rare in developed countries, with an annual incidence of 2 to 4 cases per 100,000 individuals, typically affecting older adults and individuals with underlying biliary tract disease or diabetes mellitus.

\section*{Aetiology and Pathophysiology}
The pathogenesis of liver infection varies depending on the specific infectious agent, but it generally involves cellular injury mediated either by direct cytopathic effects or by the host's immune response.
\begin{itemize}
    \item \textbf{Hepatitis A virus (HAV) and Hepatitis E virus (HEV):} Transmitted via the faecal-oral route, usually through contaminated food or water. They cause acute, self-limiting infections and do not lead to chronic disease, although HEV can cause severe, high-mortality infection in pregnant women.
    \item \textbf{Hepatitis B virus (HBV):} Transmitted through contact with infectious blood, semen, and other body fluids. It can be acquired perinatally (vertical transmission), during early childhood, or through sexual contact and needle sharing. The host immune response against infected hepatocytes is the primary driver of liver cell necrosis and fibrosis.
    \item \textbf{Hepatitis C virus (HCV):} An enveloped RNA virus transmitted primarily through blood-to-blood contact. Over 75\% to 85\% of exposed individuals fail to clear the virus acutely, resulting in chronic infection characterized by persistent hepatic inflammation and progressive fibrosis.
    \item \textbf{Pyogenic liver abscess:} A bacterial infection resulting in a localized collection of pus. Pathogens (most commonly \textit{Escherichia coli}, \textit{Klebsiella pneumoniae}, and \textit{Bacteroides} species) gain access to the liver via the biliary tract (ascending cholangitis), the portal vein (pylephlebitis secondary to appendicitis or diverticulitis), or systemic circulation.
    \item \textbf{Amoebic liver abscess:} Caused by the protozoan \textit{Entamoeba histolytica}, which migrates from the colon to the liver via the portal circulation, leading to enzymatic tissue liquefaction and abscess formation.
\end{itemize}

\section*{Clinical Features}
The clinical manifestations of liver infections range from completely asymptomatic cases to severe, life-threatening hepatic failure.
\begin{itemize}
    \item \textbf{Jaundice:} Yellowing of the skin and sclera due to impaired bilirubin excretion.
    \item \textbf{Systemic symptoms:} Fever, chills, fatigue, malaise, myalgia, and arthralgia (especially in acute viral prodromes).
    \item \textbf{Gastrointestinal features:} Nausea, vomiting, anorexia, right upper quadrant abdominal pain, and abdominal distension.
    \item \textbf{Excretory changes:} Dark, tea-coloured urine and pale, clay-coloured stools due to biliary obstruction or cholestasis.
    \item \textbf{Pruritus:} Intense, generalised itching caused by the deposition of bile salts in the skin.
    \item \textbf{Hepatomegaly:} An enlarged, tender liver felt on abdominal palpation.
\end{itemize}

\section*{Differential Diagnosis}
Liver infections must be differentiated from other inflammatory, metabolic, and obstructive causes of hepatic dysfunction.
\begin{itemize}
    \item \textbf{Drug-induced liver injury (DILI):} Hepatotoxic reactions caused by medications (such as paracetamol overdose, certain antibiotics, or herbal supplements) that mimic viral hepatitis.
    \item \textbf{Alcoholic hepatitis:} Severe liver inflammation associated with heavy, long-term alcohol consumption, presenting with similar right upper quadrant pain and jaundice.
    \item \textbf{Metabolic dysfunction-associated steatohepatitis (MASH):} A progressive form of non-alcoholic fatty liver disease that causes chronic elevation of transaminases and liver inflammation.
    \item \textbf{Autoimmune hepatitis:} A chronic inflammatory liver disease characterized by circulating autoantibodies and hypergammaglobulinaemia.
    \item \textbf{Biliary obstruction:} Cholelithiasis (gallstones), choledocholithiasis, or biliary malignancies that present with acute right upper quadrant pain, fever, and jaundice (Charcot's triad), mimicking pyogenic liver infections.
\end{itemize}

\section*{When to Seek Medical Care}
Individuals who develop unexplained fatigue, right upper quadrant pain, fever, or notice a yellowing of their skin or eyes should seek medical attention promptly.

\textbf{Contact your local emergency services for:}
\begin{itemize}
    \item Sudden onset of confusion, severe drowsiness, or a flapping tremor of the hands (signs of hepatic encephalopathy)
    \item Vomiting blood (haematemesis) or passing dark, tarry stools (melaena), indicating bleeding gastrointestinal varices
    \item High fever with shaking chills (rigors) and confusion, suggesting sepsis or a ruptured liver abscess
    \item Severe, persistent abdominal pain associated with a rigid, guard-like abdomen
    \item Sudden, significant abdominal swelling (ascites) accompanied by shortness of breath
\end{itemize}

\section*{Diagnosis and Investigations}
A combination of serology, biochemical markers, and imaging is used to diagnose the specific type of liver infection and assess the extent of liver damage.
\begin{itemize}
    \item \textbf{Liver function tests (LFTs):} Characterised by marked elevation of alanine aminotransferase (ALT) and aspartate aminotransferase (AST) (often $>1000$~U/L in acute viral infections). Bilirubin, alkaline phosphatase (ALP), and gamma-glutamyl transferase (GGT) are also elevated.
    \item \textbf{Viral serology:} Specific antibody and antigen assays, including HAV IgM, HBsAg, HBcAb IgM (to diagnose acute HBV), and anti-HCV antibodies.
    \item \textbf{Molecular testing:} Detection of HBV DNA or HCV RNA via polymerase chain reaction (PCR) to confirm active replication and guide treatment.
    \item \textbf{Imaging studies:} Abdominal ultrasound is the primary modality to evaluate liver texture, detect focal lesions (such as abscesses), and assess the biliary tree. Contrast-enhanced computed tomography (CT) or magnetic resonance imaging (MRI) is indicated for detailed evaluation of abscesses or suspected tumours.
    \item \textbf{Microbiological testing:} Blood cultures and aspirate fluid cultures (for bacterial or amoebic abscesses) to identify the causative organism and determine antibiotic sensitivity.
    \item \textbf{Non-invasive fibrosis assessment:} Transient elastography (FibroScan) or serum biomarkers (such as APRI or FIB-4) to stage chronic liver damage.
\end{itemize}

\section*{Management}
Management strategies depend entirely on the underlying pathogen and whether the infection is acute, chronic, or suppurative.
\begin{itemize}
    \item \textbf{Supportive care:} For acute, self-limiting viral infections (HAV and HEV), focusing on rest, adequate hydration, avoidance of alcohol, and discontinuation of non-essential hepatotoxic medications.
    \item \textbf{Chronic HBV therapy:} Long-term treatment with oral nucleos(t)ide analogues (such as tenofovir or entecavir) to suppress viral replication, reduce hepatic inflammation, and prevent progression to cirrhosis.
    \item \textbf{HCV cure:} A course of direct-acting antivirals (DAAs, such as sofosbuvir/velpatasvir or glecaprevir/pibrentasvir) taken for 8 to 12 weeks, which achieves viral eradication (cure) in over 95\% of patients.
    \item \textbf{Pyogenic abscess intervention:} Treatment with prolonged courses of intravenous antibiotics (such as ceftriaxone combined with metronidazole), paired with ultrasound- or CT-guided percutaneous needle aspiration or catheter drainage.
    \item \textbf{Amoebic abscess intervention:} High-dose oral metronidazole or tinidazole, followed by a luminal amoebicide (such as paromomycin) to eradicate the intestinal reservoir.
\end{itemize}

\section*{Complications}
\begin{itemize}
    \item \textbf{Hepatic cirrhosis:} Chronic inflammation leading to diffuse nodular regeneration and extensive scarring of the liver tissue.
    \item \textbf{Hepatocellular carcinoma (HCC):} Primary liver cancer arising in the setting of chronic HBV or HCV infection and cirrhosis.
    \item \textbf{Portal hypertension:} Increased pressure in the portal venous system, leading to splenomegaly, ascites, and life-threatening oesophageal varices.
    \item \textbf{Hepatic encephalopathy:} Cognitive decline, neuromuscular impairment, and coma due to the accumulation of neurotoxins (like ammonia) in the blood.
    \item \textbf{Acute liver failure:} Rapid loss of hepatic function leading to coagulopathy and encephalopathy, primarily seen in severe acute HBV, HEV, or drug-toxicity cases.
    \item \textbf{Septic shock:} Systemic dissemination of bacteria from a pyogenic liver abscess, carrying a high mortality rate.
\end{itemize}

\section*{Prognosis and Outcomes}
The prognosis of liver infections is highly variable and depends on the specific cause and early initiation of treatment. Acute hepatitis A and E generally carry an excellent prognosis, with complete clinical recovery and no long-term sequelae in the vast majority of patients. Chronic hepatitis B cannot currently be cured, but suppression of viral replication with daily medication dramatically reduces the risk of cirrhosis and liver cancer. Chronic hepatitis C is now highly curable, and successful treatment halts the progression of liver damage, though patients with pre-existing cirrhosis require ongoing surveillance for hepatocellular carcinoma. Pyogenic liver abscesses carry a good prognosis with timely drainage and targeted antibiotics, but delays in diagnosis can lead to overwhelming sepsis and death.

\section*{Prevention and Self-Care}
\begin{itemize}
    \item \textbf{Vaccination:} Complete the full vaccination series for hepatitis A and hepatitis B, which is highly recommended for all infants, travellers, and high-risk individuals in many countries.
    \item \textbf{Safe injecting practices:} Never share needles, syringes, or other drug preparation equipment. Utilise supervised injecting facilities and needle syringe programmes.
    \item \textbf{Safe sexual practices:} Use barrier contraception (condoms) consistently to reduce the risk of sexually transmitted hepatitis B.
    \item \textbf{Hygiene and sanitation:} Practise thorough hand hygiene, particularly after using the bathroom and before preparing food. Drink bottled water and avoid raw shellfish when travelling to regions with poor sanitation.
    \item \textbf{Avoid sharing personal care items:} Do not share toothbrushes, razors, or nail clippers, which can carry microscopic amounts of blood.
    \item \textbf{Ensure sterile procedures:} Only obtain tattoos, piercings, or medical procedures from licensed studios and clinics that use single-use, sterile equipment.
\end{itemize}

\section*{Sources and Further Reading}
The following organisations provide authoritative information and clinical guidelines on liver infections and viral hepatitis.
\begin{itemize}
    \item Gastroenterological Society of Australia (GESA). \textit{Clinical guidelines for hepatitis B and C management.} https://www.gesa.org.au/
    \item Hepatitis Australia. \textit{Support, advocacy, and education resources for viral hepatitis.} https://www.hepatitisaustralia.com/
    \item World Health Organization. \textit{Global hepatitis report and prevention strategies.} https://www.who.int/
    \item Healthdirect Australia. \textit{Hepatitis and liver infections.} https://www.healthdirect.gov.au/hepatitis
    \item Mayo Clinic. \textit{Hepatitis: symptoms, causes, and diagnosis.} https://www.mayoclinic.org/
\end{itemize}
